% a06_BALtorus_d2: part of Step 3 of the proof of thm:a06:T1 (input from sec/a06_BALtorus_c.tex).
% Group (D), first half: the counterterm moves of Ba3 pp. 411-422, the vacuum functionals of a datum and their exact algebra.

\subsubsection*{Step 3(D$_1$): the counterterm moves of \cite[pp.~411--422]{Ba3}, exact algebra}

Two statements of \cite{Ba3} on the vacuum part of the effective action carry no printed argument: that the field-independent part generated after $k$ steps on the domains of the step, measured against the localised counterterms $E_k(\cdot,\Omega_1)$ of \cite[(1.30)]{Ba3}, is $O(t_k)$ per unit block (``obvious consequences of Proposition~1'', \cite[pp.~421--422]{Ba3}), and that the domain and localisation changes of a step, with the counterterm move of \cite[p.~422]{Ba3}, cost only collar terms (``even simpler and we omit them''). We write them as the statements $(\mathrm{VAC}^\flat)_k$ and $(\mathrm{DOM})_k$ (\cref{thm:a06d:V12,thm:a06d:V14}) and prove them from Proposition~1. This part sets up the objects and their exact identities; the bounds are in Step~3(D$_2$).

\begin{definition}[datum at scale $k$]\label{def:a06d:V1}
A datum is $\mathfrak D=(\Omega,\Omega_1,g,c)$: $\Omega$ a union of big blocks of $T^{(k)}_1$ (the domain hypothesis of \cite[Props.~2.1--2.3]{Ba1}, \cite[p.~411]{Ba3}), on whose $\eta$-points the fine scalar field lives with the Neumann covariant Laplacian; $\Omega_1\subset\Omega$ a union of unit blocks; $g$ smooth with $\supp g\subset\Omega_1$, $g=1$ where $\dist(\cdot,\partial\Omega_1)\ge M$ and $\norm g_{C^2}\le C_g$ (\cite[p.~412]{Ba3} has $C_g=1$; a fixed $C_g$ is allowed, as for the partitions of unity of \cite[p.~420]{Ba3}); $c(x)$ the coefficient of the vertices \cite[(1.7)]{Ba3}, namely \cite[(1.29)]{Ba3} built with $g$ at the scale of $\mathfrak D$, $c=\dm_{\mathrm{fin}}t_k^2+\dm_{K,k}t_k^2+\dm_k(se(t_k)g,\lambda'\lambda(t_k),\Omega_1,x)$. The depth is $\rho(\mathfrak D)\defeq\dist(\Omega_1,\Omega^c)$ in scale-$k$ units. Ba{\l}aban's step-$k$ datum is $\mathfrak D_k\defeq(B^k(\Lambda^{(k-1)'}_2),B^k(\Lambda^{(k-1)'}_7),g_k,c_k)$; torus data are $\mathfrak D^T_{W,g}\defeq(T,W,g,c_{W,g})$, and the plain torus datum $\mathfrak D^T$ has $W=T$, $g=1$, $c=\dm$.
\end{definition}

\begin{definition}[the vacuum functionals of a datum]\label{def:a06d:V2}
$\Xi_k[\mathfrak D](A,\phi)$ is the nested integral inside the logarithm of \cite[(1.4)]{Ba3}, with $(\Omega,\Omega_1,g_k)$ and the coefficient replaced by those of $\mathfrak D$, $e'$ replaced by $s$, and the counterterm factor $e^{-E_1}$ omitted (the vector covariances are the torus $C^{(j)}$ of \cite[(1.4)]{Ba3}). $\mathrm{Gen}_k[\mathfrak D]\defeq T_{[1,\bar n]}\log[\Xi_k[\mathfrak D](0,0)/\Xi_k[\mathfrak D](0,0)|_{s=\lambda'=0}]$ (the generated vacuum constants); $F_k[\mathfrak D](A,\phi)\defeq T_{[0,\bar n]}\log[\Xi_k[\mathfrak D](A,\phi)/\Xi_k[\mathfrak D](0,0)]$ (the field-dependent part, including its degree-$0$ Gaussian part); $\mathrm{Ct}_k[\mathfrak D]\defeq E_k(se(t_k)g,\lambda'\lambda(t_k),\Omega_1)$ of \cite[(1.30)]{Ba3}, under \cref{conv:a06d:LOC}; $\mathcal R_k[\mathfrak D]\defeq\mathrm{Gen}_k[\mathfrak D]-\mathrm{Ct}_k[\mathfrak D]$.
\end{definition}

By \cite[(1.4)]{Ba3}, $E_k(\cdot,\Omega,0,0)=E_1(e',g_k,\lambda',\Omega_1)-\log\Xi_k[\mathfrak D_k](0,0)$, and by \cite[(1.30)]{Ba3}, $E_1(e',g_k,\lambda',\Omega_1)=\Efin+E_{K,k}+\mathrm{Ct}_k[\mathfrak D_k]$; so the field-independent part of the polynomial $\mathcal P^{(k)}(\Omega_1,\cdot)$ of \cite[(1.5)]{Ba3}, in the logarithmic convention, is $\mathcal R_k[\mathfrak D_k]-\Efin-E_{K,k}$, with $\Efin+E_{K,k}$ a global front factor.

\begin{convention}[how the localised counterterms are built]\label{conv:a06d:LOC}
\textup{(a)} Every vertex \cite[(1.7)]{Ba3} carries the coefficient $c$ of \cite[(1.29)]{Ba3} as a number, not opened, both in the expansion of \cite[(1.4)]{Ba3} and in the graphs of $\mathrm{Ct}_k$. \textup{(b)} $\mathrm{Ct}_k[\mathfrak D]$ is the sum of the bare vacuum graphs (vertices \cite[(1.6), (1.7), (1.8), (1.10)]{Ba3} with $n'=0$, $\tilde B=0$) of orders $1\le\alpha+\beta\le\bar n$, with every line of index $<k$: scalar lines $G_k(T,0)$ (the rescaled $G^\eps_k(0)$ of \cite[(1.25)--(1.28)]{Ba3}) and vector lines $G^{\mathrm{vec}}_k(T)$, the vertices localised exactly as the generated ones of \cite[(1.4)]{Ba3}: \cite[(1.6)--(1.7)]{Ba3} summed over $\Omega_1$, $g$ on each vector leg. \textup{(c)} In $\dm_k(se(t_k)g,\dots)$ the factor $g$ sits on the vector legs of the coefficient's graphs, so a divergent self-energy subgraph of a generated graph and its \cite[(1.7)]{Ba3} partner (\cite[p.~420]{Ba3}: ``the same graph $G$ but with both external legs localized in $x$'') carry $g$ on corresponding legs.
\end{convention}

(b) and (c) are printed: \cite[(1.31), p.~419]{Ba3} displays the localised second-order photon counterterm with torus lines and $g_k$ on both vector legs, and \cite[p.~419]{Ba3} reads ``localize the vertices (1.6) and (1.7) restricting summations to the set $\Omega_1$. The other vertices are localized by functions $g_k$ multiplying each leg of the vector field''. (a) is forced: if the \cite[(1.7)]{Ba3} vertices of $\mathrm{Ct}$ carried the full torus $\dm$, $\mathrm{Ct}$ would contain on $\{g<1\}$ self-energy subgraphs with legs $gg'$ paired with an unlocalised coefficient, an uncancelled $O(e(t)^2\eta^{-1})$ per site; if the coefficients were opened into lines of index $<k$, $\mathrm{Gen}-\mathrm{Ct}$ would contain $(\dm_{\mathrm{fin}}t^2+\dm_{K,k}t^2)G_k(x,x)=O(t\eta^{-1})$ per site. On the torus datum, $\Efin+E_{K,k}+\mathrm{Ct}_k[\mathfrak D^T]=E_1$. (Ba{\l}aban remarks on \cite[p.~418]{Ba3} that it ``is sufficient'' to keep the orders $2\le\alpha+2\beta\le6$ in $E_1$; under that restricted $E_1$ the convergent bare graphs of higher order are extra standalone renormalised families in $\mathcal R_k$, bounded by \cref{lem:a06d:V9}(a) below, and nothing else changes.)

\begin{definition}[graphs, families, localisations]\label{def:a06d:V4}
Expanding $\log\Xi_k[\mathfrak D](0,0)$ in $(s,\lambda')$ gives connected vacuum graphs \cite[pp.~414--415]{Ba3}, and $\mathrm{Gen}_k=\mathrm{Gen}^{\mathrm{bare}}_k+\mathfrak B_k$: bare graphs (all vertices from \cite[(1.6), (1.7), (1.8), (1.10)]{Ba3}, from the expansion of $-\Delta^\eta_{sgA'}$ and of $V$) and contour graphs (at least one product \cite[(1.18)]{Ba3} of the form $(1.13)\cdot(1.14)$ or $(1.14)\cdot(1.14)$, from the expansion of $a_k|Q_k(sgA')\phi'|^2-a_k|Q_k(0)\phi'|^2$; \cite[(1.12)]{Ba3} does not occur since the block field is $0$). Scalar lines are $G_k(\Omega,0)$ (at $\phi=0$ the Gaussian part is $-\frac12\langle\phi',(-\Delta^\eta+m^2t_k^2+a_kP_k(0))\phi'\rangle$ on $\Omega$); vector lines are $G^{\mathrm{vec}}_k(T)$, the covariance of $A'$ \cite[p.~414]{Ba3}. Localisation follows \cite[p.~420]{Ba3}, the families are the $G_{\mathrm{ren}}$ of \cite[p.~420]{Ba3}, $X(\Lambda)$ is the sum of the terms of $X$ with localisation $\Lambda=\{\square(v)\}$, and $d(\Lambda)$ the tree length of \cite[(1.33)]{Ba3}.
\end{definition}

\begin{lemma}[the generated bare part and the counterterm differ only in their lines]\label{lem:a06d:V5}
Under \cref{conv:a06d:LOC}, $\mathrm{Gen}^{\mathrm{bare}}_k[\mathfrak D]$ and $\mathrm{Ct}_k[\mathfrak D]$ are sums over the same graphs, with the same localisations and vertex coefficients, with lines $(G_k(\Omega,0),G^{\mathrm{vec}}_k(T))$ and $(G_k(T,0),G^{\mathrm{vec}}_k(T))$ respectively. Hence $\mathcal R_k[\mathfrak D]=\mathfrak B_k[\mathfrak D]+\Delta_k[\mathfrak D]$ with
\[
\Delta_k[\mathfrak D]\defeq\sum_G\sum_\ell G\big[G_k(T,0)\text{ on the scalar lines before }\ell;\ \delta_k\text{ on }\ell;\ G_k(\Omega,0)\text{ after }\ell\big],\qquad\delta_k\defeq G_k(\Omega,0)-G_k(T,0)=-\delta G_k(T,\Omega,0),
\]
the scalar lines of each $G$ ordered so that the lines of each closed scalar loop are consecutive.
\end{lemma}

\begin{proof}
At zero block fields the bare vertex functions of \cite[(1.4)]{Ba3} come from the Taylor expansion of $U(sgA'_b)=\exp(qs\eta g(b)A'_b)$ inside $-\Delta^\eta_{sgA'}$ and from $V$: they are \cite[(1.8), (1.10)]{Ba3} with $n'=0$, $\tilde B=0$, together with \cite[(1.6), (1.7)]{Ba3}; the bare vacuum graphs of $S^\eps$ \cite[p.~416]{Ba3} come from the same expansion of the lattice covariant Laplacian, rescaled, and \cref{conv:a06d:LOC}(a), (b) give them the same localisation and coefficients. The generated scalar lines are $G_k(\Omega,0)$; the counterterm lines have index $<k$, and by \cite[(1.27)]{Ba3} the index-$<k$ part of $G^\eps_K(0)$ is $G^\eps_k(0)$, which rescales to $G_k(T,0)$; the vector lines agree. Multilinearity in the scalar lines, $\prod_\ell a_\ell-\prod_\ell b_\ell=\sum_\ell(\prod_{\ell'<\ell}b_{\ell'})(a_\ell-b_\ell)(\prod_{\ell'>\ell}a_{\ell'})$, gives $\Delta_k$. $S^\eps$ has no block kernel, so no vertex \cite[(1.13)--(1.15)]{Ba3} occurs in $\mathrm{Ct}_k$, and the contour graphs, which have no counterterm, stay in $\mathfrak B_k$.
\end{proof}

The single terms of $\Delta_k$ are never bounded one by one: the vacuum graph made of one \cite[(1.7)]{Ba3} vertex and one scalar loop gives $c(x)\delta_k(x,x)$ with $c\sim e(t)^2\eta^{-1}$, cancelled only together with the graph carrying the self-energy $S$ with $\delta_k$ on the closing line, $[S-c]\delta_k=[S-S(T)]\delta_k+[S(T)-c]\delta_k$. So $\Delta_k$, and every difference below, is used only through renormalised families whose members carry the telescoping order induced from $G$ on $G/S$ (\cref{lem:a06d:K2}); with the ordering of \cref{lem:a06d:V5} every closed scalar loop not containing $\ell$ carries uniform full lines, all $G_k(T,0)$ or all $G_k(\Omega,0)$, and no closed loop is ever decomposed by fixed scales or orderings. In the measure convention of \cite[(2.1)]{Ba3} a lone tadpole has $D=-1$ and Proposition~1 never bounds it alone.

\begin{lemma}[exact vacuum staging on a fixed datum]\label{lem:a06d:V6}
Fix a datum $\mathfrak D$ at scale $k$ and let $\mathfrak D\!\uparrow$ be the same datum read at scale $k+1$ (same fine domains, vertex functions and coefficient $c$). Let $d\mu_k$ be the Gaussian law of the step-$k$ fluctuations at zero $(k+1)$-block fields with the composing covariances $C^{(k)}$ (vector, torus) and $C^{(k)}(\Omega)\defeq(\Delta^{(k)}(\Omega,0)+\alpha Q^*Q)^{-1}$ (scalar, \cref{lem:a06d:B6} with $H$ the Neumann operator on $\Omega$, $B=0$). For a datum with localisation $g$ let $\iota^g_kf$ be the Taylor series of $\tau\mapsto f(\tau gA'^{(k),\eps}+eA^{(k+1),\eps})$ at $\tau=0$ with $\tau\to s$ ($\iota^g_k=\iota_k$ on the torus; \cref{lem:a06d:B4,lem:a06d:B5}(a) hold verbatim for it, line integrals of $gA'$ being additive). Put
\[
\mathrm{out}_k[\mathfrak D]\defeq T_{[1,\bar n]}\log\int d\mu_k\exp\big(\iota^g_kF_k[\mathfrak D]-(\iota^g_kF_k)_0+\mathcal K_k\big),\quad(\iota^g_kF_k)_0\defeq F_k|_{s=\lambda'=0,A=0}=-\tfrac12\langle\phi,\Delta^{(k)}(\Omega,0)\phi\rangle,\quad\mathcal K_k\defeq-\tfrac12\alpha\big(|Q(sgA'^{(k)})\phi|^2-|Q(0)\phi|^2\big),
\]
the vacuum constant generated by step $k+1$. Then: \textup{(a)} $\mathrm{Gen}_{k+1}[\mathfrak D\!\uparrow]=\mathrm{Gen}_k[\mathfrak D]+\mathrm{out}_k[\mathfrak D]$; \textup{(b)} the lines of $\mathrm{Gen}_{k+1}[\mathfrak D\!\uparrow]$ are $G_{k+1}(\Omega,0)=G_k(\Omega,0)+\beta_k^2G_k(\Omega,0)R_k^*C^{(k)}(\Omega)R_kG_k(\Omega,0)$ and $G^{\mathrm{vec}}_{k+1}=G^{\mathrm{vec}}_k+(\text{lifted }C^{(k)})$, so $\mathrm{out}_k[\mathfrak D]$ is the part of $\mathrm{Gen}_{k+1}[\mathfrak D\!\uparrow]$ with at least one line of index $k$; \textup{(c)} with $\mathcal V_k[\mathfrak D]\defeq\mathrm{Ct}_{k+1}[\mathfrak D\!\uparrow]-\mathrm{Ct}_k[\mathfrak D]$ (the bare vacuum graphs with all lines of index $\le k$ and at least one of index $k$, localised and coefficiented as in $\mathfrak D$), exactly
\[
\mathcal R_{k+1}[\mathfrak D\!\uparrow]=\mathcal R_k[\mathfrak D]+\mathrm{out}_k[\mathfrak D]-\mathcal V_k[\mathfrak D].
\]
\end{lemma}

\begin{proof}
Fubini on the nested integral: the $(k+1)$-step composite kernel on $\Omega$ factors into the step-$(k+1)$ kernels and the $k$-step composite (\cref{lem:a06d:B3,lem:a06d:B4,lem:a06d:B5}(a), all valid on $\Omega$), the gauge part being \cref{lem:a06d:B6}(b). So $\Xi_{k+1}[\mathfrak D\!\uparrow](0,0)=N\int d\mu_k\exp(\iota^g_kF_k[\mathfrak D]-(\iota^g_kF_k)_0+\mathcal K_k)\,\Xi_k[\mathfrak D](0,0)$ up to terms of degree $>\bar n$: only the zero-background degree-$0$ part $-\frac12\langle\phi,\Delta^{(k)}(\Omega,0)\phi\rangle$ and the term $\frac12\alpha|Q(0)\phi|^2$ combine with the step kernels into the Gaussian weight of $d\mu_k$, and every other term of $\iota^g_kF_k$, and $\mathcal K_k$, has degree $\ge1$ and is kept. (The degree-$0$ part of $F_k$ depends on the background $A$ through the scalar fluctuation determinants, Ba{\l}aban's $Z^{(j)}$ factors \cite[(2.52), (2.92), (2.105)]{Ba2}, and through $-\frac12\langle\phi,\Delta^{(k)}(A)\phi\rangle$; $\iota^g_k$ turns that dependence into terms of degree $\ge1$, which belong to $\mathrm{out}_k$, as on \cite[p.~576]{Ba2}: ``in the density $\varrho'^{(k)}$ only the factors $Z^{(j)}$ depend on this field''. $\mathcal K_k$ is the fluctuation dependence of the new scalar kernel at $\psi=0$; it gives the \cite[(1.12)]{Ba3}-type contour vertices on the new outer segment $\Gamma_{z,x_k}$, and without it (a) is false.) Take logarithms; constants pass through, and \cref{lem:a06d:B2} truncates: this is (a). (b) is \cref{lem:a06d:B6}(c), with $H$ the Neumann operator on $\Omega$ for the scalar and on the torus for the vector; expanding multilinearly in the lines identifies $\mathrm{out}_k$ with the terms having at least one index-$k$ line. (c): $\mathcal R_{k+1}-\mathcal R_k=(\mathrm{Gen}_{k+1}-\mathrm{Gen}_k)-(\mathrm{Ct}_{k+1}-\mathrm{Ct}_k)$, the coefficient $c$ being the same in $\mathrm{Gen}$ and $\mathrm{Ct}$ and held fixed between $k$ and $k+1$.
\end{proof}

On the torus datum $\Delta=0$ and (c) reads $\mathrm{out}_k=\mathcal V_k+\mathfrak B_{k+1}-\mathfrak B_k$; on a general datum (c) holds with the $\Delta$ terms, and a term of $\mathcal R_k$ is never bounded in absolute value except where a datum changes.

\begin{lemma}[comparison of data]\label{lem:a06d:V7}
Let $\mathfrak D,\mathfrak D'$ be data at the same scale. \textup{(a)} (Propagator domain.) If $\mathfrak D'$ is $\mathfrak D$ with $\Omega$ replaced by an admissible $\Omega'\subset\Omega$ and the same vertex datum, then $\mathcal R_k[\mathfrak D]-\mathcal R_k[\mathfrak D']$ is a sum of families of the two forms of \cref{lem:a06d:V5} (contour or $\Delta$), each with one designated line $\delta G_k(\Omega,\Omega',0)$ (in Ba{\l}aban's list, \cite[p.~414]{Ba3}); likewise $F_k[\mathfrak D]-F_k[\mathfrak D']$ is a sum of families with one designated $\delta G_k(\Omega,\Omega',\cdot)$ line or one designated $\delta$ in the external-field map $a_kG_k(\cdot)Q_k^*$ \cite[p.~415]{Ba3}. \textup{(b)} (Vertex datum.) If $\mathfrak D'$ is $\mathfrak D$ with $(\Omega_1,g,c)$ replaced by $(\Omega'_1,g',c')$, then $\mathcal R_k[\mathfrak D]-\mathcal R_k[\mathfrak D']$ (likewise $\mathrm{out}_k$, $F_k$) is a sum of families of the forms of \cref{lem:a06d:V5} with one designated vertex of one of three kinds: (b1) a factor $(g-g')$ on one vector leg, with the same change on the corresponding leg of the \cite[(1.7)]{Ba3} partner's coefficient graph (\cref{conv:a06d:LOC}(c)); (b2) a \cite[(1.6)]{Ba3} or \cite[(1.7)]{Ba3} vertex in $\Omega_1\triangle\Omega'_1$; (b3) a \cite[(1.7)]{Ba3} vertex carrying a change of the $g$-independent finite part of $c$ ($\dm_{\mathrm{fin}}$, $\dm_{K,\cdot}$, or a coefficient graph with an index-$k$ line, all regular; $|\text{change}|\le C_Tt_k(1+\ln t_k^{-1})^5$ per site by \cref{thm:a06e:T0}, and $\dm_{K,k}(L^k\eps)^2=O(L^k\eps)$ by \cref{thm:a06e:DM}, as asserted on \cite[p.~418]{Ba3}). The designated factor is supported in $Z(\mathfrak D,\mathfrak D')\defeq(\Omega_1\triangle\Omega'_1)\cup\supp(g-g')\cup\supp(c-c')$ up to exponentially decaying tails of the regular coefficient graphs. \textup{(c)} (Step covariances.) For two choices $C,C'$ of the step-$k$ covariances, $\mathrm{out}_k[\mathfrak D;C]-\mathrm{out}_k[\mathfrak D;C']$ is a sum of step-$k$ cumulant families with one designated contraction by $\delta C\defeq C-C'$, lifted by $\beta_kG_kR_k^*$ (\cref{lem:a06d:B6}(b)).
\end{lemma}

\begin{proof}
Telescope multilinearly: in the lines for (a), in the vertex multipliers for (b), in the contraction kernels for (c). In (b) the telescoping is symmetric, $g'=g+(g'-g)$ and $\mathbf 1_{\Omega'_1}=\mathbf 1_{\Omega_1}+(\mathbf 1_{\Omega'_1}-\mathbf 1_{\Omega_1})$ on all legs and sites at once, applied identically to $\mathrm{Gen}$, $\mathrm{Ct}$ and the coefficient graphs; every term then carries at least one factor $(g'-g)$ or a site difference, the members of a renormalised pair carry the same set of changed legs and sites (\cref{conv:a06d:LOC}(c)), and the number of terms is $\le2^{\text{legs}+\text{sites}}\le C(\bar n)$. (A lone \cite[(1.7)]{Ba3} vertex is divergent and the \cite[(1.6)]{Ba3} tadpole needs its partner, so the kinds are not bounded vertex by vertex.) By \cref{conv:a06d:LOC} $\mathrm{Gen}$ and $\mathrm{Ct}$ carry the same vertices, so each difference term again has the form $\mathfrak B+\Delta$ of \cref{lem:a06d:V5}: no bare vacuum graph with an uncancelled overall divergence survives; by (c) of the convention the change on a self-energy subgraph and on its partner are made on corresponding legs, so renormalised families stay renormalised. A coefficient change splits into the $g$-localised part of \cite[(1.29)]{Ba3}, $\dm_k(\cdot g)$, which changes with $g$ and is paired with the generated subgraphs (kind (b1)), and the $g$-independent finite parts, which change only with the scale index ($\dm_{K,k}\to\dm_{K,k+1}$), the moved piece being an index-$k$ slice, a regular graph (kind (b3)). Here $(g'-g)$ is itself a multiplier of the class ($C^2$-norm $\le2C_g$) supported in $Z$.
\end{proof}
